\section{Exercício prático: servidor Nginx}

Neste exercício, o objetivo é subir um servidor Nginx, alterar a página inicial via terminal e monitorar os logs em tempo real.

\subsection{Subindo o servidor Nginx}

\begin{verbatim}
docker run -d -p 8080:80 --name meu-nginx nginx
\end{verbatim}

Esse comando inicia um container Nginx em segundo plano, mapeando a porta 8080 do host para a porta 80 do container. Acessando \texttt{http://localhost:8080} no navegador, deve aparecer a página padrão do Nginx.

\subsection{Alterando a página inicial via terminal}

Usando \texttt{docker exec}, é possível acessar o shell do container e modificar o arquivo HTML da página padrão:

\begin{verbatim}
docker exec -it meu-nginx bash

# Dentro do container:
echo '<h1>Minha pagina Docker!</h1>' > /usr/share/nginx/html/index.html
exit
\end{verbatim}

Ao recarregar \texttt{http://localhost:8080} no navegador, a nova página aparecerá. O diretório \texttt{/usr/share/nginx/html/} é o diretório padrão onde o Nginx serve arquivos estáticos.

\subsection{Monitorando logs em tempo real}

\begin{verbatim}
docker logs -f meu-nginx
\end{verbatim}

A flag \texttt{-f} segue os logs em tempo real. Cada vez que a página for acessada no navegador, uma nova linha de log de acesso aparecerá no terminal. Para sair, basta pressionar \texttt{Ctrl+C}.

Para ver apenas as últimas linhas:

\begin{verbatim}
docker logs --tail 10 meu-nginx
\end{verbatim}

\subsection{Limpeza}

\begin{verbatim}
docker rm -f meu-nginx
\end{verbatim}
